\begin{activity}{Buried Sphere --- Width to Depth}

\section*{Purpose}

This activity derives and applies a simple depth-estimation rule for gravity anomalies using the buried sphere model. By working through the mathematics of the full-width at half-maximum, you will connect the observable width of an anomaly to the depth of its source, independent of the source mass --- a result that illustrates both the utility and the non-uniqueness of gravity interpretation.

\section*{Learning Outcomes}

By the end of this activity, you should be able to:
\begin{itemize}
\item Relate observable characteristics of a gravity anomaly to source geometry and depth.
\item Derive a simple relationship between anomaly width and source depth from the buried sphere model.
\item Explain why some source parameters influence anomaly amplitude whereas others influence anomaly shape.
\item Describe sources of non-uniqueness in gravity interpretation and suggest ways to reduce them.
\item Use simple analytical results to guide preliminary geological interpretation prior to inversion.
\end{itemize}

\section*{Introduction}
The gravity anomaly due to a buried sphere is given by:
\[
  g_z(x) = \tfrac{4}{3}\pi a^3 G\, \Delta\rho 
  \frac{z_c}{(x^2+z_c^2)^{3/2}}\, ,
\]
where $a$ is the radius of the sphere, $\Delta rho$ is its density anomaly, and $z_c$ is the depth to the center of the sphere.  The coordinate $x$ is the position along the profile at Earth's surface.

Last week, we saw this relationship leads to non-uniqueness because $a^3 \Delta \rho$ are inseperable.  Now, however, we want to show that one can produce a simple relationship between the depth of the buried sphere and the width of its anomaly, independent of the the mass.  Such a relationship is useful for estimating the depth to structures that produce peaks in the gravity field.  Such a relationship is useful for producing back of the envelope estimates in the field or from a map before any formal modeling is done.  Such relationships are the prelude concepts that lead to wavelength filtering that we will see in Week 7.

\begin{figure}[h]
  \centering
  \begin{tikzpicture}[scale=2.0]
    % parameter
    \def\z{1.0}

    % axes
    \draw[->] (-3,0) -- (3.2,0) node[right] {$x$};
    \draw[->] (0,-0.1) -- (0,1.3) node[above] {$g(x)$};

    % function
    \draw[thick,black]
        plot[domain=-3:3,samples=200]
        (\x,{1/((\x*\x+\z*\z)^1.5)});

    % half-width indicator
    \draw[dashed] (-0.77,0.5) -- (-0.77,-0.1);
    \draw[dashed] ( 0.77,0.5) -- ( 0.77,-0.1);

    \draw[<->] (-0.77,0.5) -- (0.77,0.5);

    % labels
    \node[below] at (-0.77,-0.2) {$-x_{1/2}$};
    \node[below] at ( 0.77,-0.2) {$x_{1/2}$};

    \node[above] at (0.1,0.5) {$w$};

    \node[right] at (0,1.1) {$g(0)$};
    \node[right] at (0.77,0.6) {$g(x_{1/2})$};
  \end{tikzpicture}
  \caption{Gravity anomaly due to a buried sphere.}
  \label{fig:gravity_sphere}
\end{figure}

\subsection*{Tasks}
To develop the depth to width relationship:
\begin{questions}
  \question{Peak Anomaly}
  Show that $g_z(0)$ is proportional to $\Delta\rho\,a^3/z_c^2$.
  \answerbox

  \question{Width to Depth Relationship}
  Show that the depth
  \[
    z_c \approx 0.65\, w
  \]
  where \(w\) is the full-width half-maximum.
  
  To do so, solve for $x_{1/2}$ starting from $g_z(x_{1/2}) = \tfrac{1}{2} g_z(0)$.  Hint: this is on the right track,
  \[
    \frac{z_c}{x_{1/2}} = (2^{2/3} - 1)^{-1/2}.
  \]
  \answerbox[320pt]

  \clearpage
  \question{Interpretation}
    \begin{enumerate}[label=\alph*]
      \item Why does depth affect width but density contrast does not?  
      \answerbox{36pt}

      \item How do the radius and density contrast result create non-uniqueness for inversion and how we ensure uniqueness?
      \answerbox{36pt}

      \item Why does a shallow lenticular body (e.g., laccolith or lopolith) could mimic a deep sphere's anomaly?
      \answerbox{36pt}
    \end{enumerate}

  \question{When Depth-Wavelength Breaks Down}
    Explain why sedimentary basins can give broad anomalies while being shallow features.
    \answerbox

    Postulate how sedimentary basins can then be distinguished from deep features.
    \answerbox

\end{questions}

\section*{Key Ideas}

\textbf{Different source parameters influence different characteristics of an anomaly.}
Changes in density contrast primarily affect anomaly amplitude, changes in source geometry affect anomaly shape, and changes in depth affect anomaly width (or wavelength).

\textbf{Not all source properties are equally well constrained.}
Several combinations of geometry and density contrast can produce similar gravity anomalies. Understanding which parts of an anomaly are controlled by which parameters helps identify sources of non-uniqueness.

\textbf{Depth and wavelength are fundamentally related.}
Shallow sources tend to produce narrow, short-wavelength anomalies, whereas deep sources produce broader, long-wavelength anomalies. This relationship underpins many depth-estimation and filtering techniques used in potential-field geophysics.

\end{activity}